% Options for packages loaded elsewhere
% Options for packages loaded elsewhere
\PassOptionsToPackage{unicode}{hyperref}
\PassOptionsToPackage{hyphens}{url}
\PassOptionsToPackage{dvipsnames,svgnames,x11names}{xcolor}
%
\documentclass[
  french,
  10pt,
  a4paper,
  oneside]{article}
\usepackage{xcolor}
\usepackage[top=16mm,bottom=16mm,left=16mm,right=16mm,headheight=12pt,headsep=4mm,footskip=8mm,includeheadfoot,heightrounded]{geometry}
\usepackage{amsmath,amssymb}
\setcounter{secnumdepth}{-\maxdimen} % remove section numbering
\usepackage{iftex}
\ifPDFTeX
  \usepackage[T1]{fontenc}
  \usepackage[utf8]{inputenc}
  \usepackage{textcomp} % provide euro and other symbols
\else % if luatex or xetex
  \usepackage{unicode-math} % this also loads fontspec
  \defaultfontfeatures{Scale=MatchLowercase}
  \defaultfontfeatures[\rmfamily]{Ligatures=TeX,Scale=1}
\fi
\usepackage{lmodern}
\ifPDFTeX\else
  % xetex/luatex font selection
\fi
% Use upquote if available, for straight quotes in verbatim environments
\IfFileExists{upquote.sty}{\usepackage{upquote}}{}
\IfFileExists{microtype.sty}{% use microtype if available
  \usepackage[]{microtype}
  \UseMicrotypeSet[protrusion]{basicmath} % disable protrusion for tt fonts
}{}
\usepackage{setspace}
\makeatletter
\@ifundefined{KOMAClassName}{% if non-KOMA class
  \IfFileExists{parskip.sty}{%
    \usepackage{parskip}
  }{% else
    \setlength{\parindent}{0pt}
    \setlength{\parskip}{6pt plus 2pt minus 1pt}}
}{% if KOMA class
  \KOMAoptions{parskip=half}}
\makeatother
% Make \paragraph and \subparagraph free-standing
\makeatletter
\ifx\paragraph\undefined\else
  \let\oldparagraph\paragraph
  \renewcommand{\paragraph}{
    \@ifstar
      \xxxParagraphStar
      \xxxParagraphNoStar
  }
  \newcommand{\xxxParagraphStar}[1]{\oldparagraph*{#1}\mbox{}}
  \newcommand{\xxxParagraphNoStar}[1]{\oldparagraph{#1}\mbox{}}
\fi
\ifx\subparagraph\undefined\else
  \let\oldsubparagraph\subparagraph
  \renewcommand{\subparagraph}{
    \@ifstar
      \xxxSubParagraphStar
      \xxxSubParagraphNoStar
  }
  \newcommand{\xxxSubParagraphStar}[1]{\oldsubparagraph*{#1}\mbox{}}
  \newcommand{\xxxSubParagraphNoStar}[1]{\oldsubparagraph{#1}\mbox{}}
\fi
\makeatother


\usepackage{longtable,booktabs,array}
\newcounter{none} % for unnumbered tables
\usepackage{calc} % for calculating minipage widths
% Correct order of tables after \paragraph or \subparagraph
\usepackage{etoolbox}
\makeatletter
\patchcmd\longtable{\par}{\if@noskipsec\mbox{}\fi\par}{}{}
\makeatother
% Allow footnotes in longtable head/foot
\IfFileExists{footnotehyper.sty}{\usepackage{footnotehyper}}{\usepackage{footnote}}
\makesavenoteenv{longtable}
\usepackage{graphicx}
\makeatletter
\newsavebox\pandoc@box
\newcommand*\pandocbounded[1]{% scales image to fit in text height/width
  \sbox\pandoc@box{#1}%
  \Gscale@div\@tempa{\textheight}{\dimexpr\ht\pandoc@box+\dp\pandoc@box\relax}%
  \Gscale@div\@tempb{\linewidth}{\wd\pandoc@box}%
  \ifdim\@tempb\p@<\@tempa\p@\let\@tempa\@tempb\fi% select the smaller of both
  \ifdim\@tempa\p@<\p@\scalebox{\@tempa}{\usebox\pandoc@box}%
  \else\usebox{\pandoc@box}%
  \fi%
}
% Set default figure placement to htbp
\def\fps@figure{htbp}
\makeatother



\ifLuaTeX
\usepackage[bidi=basic,shorthands=off]{babel}
\else
\usepackage[bidi=default,shorthands=off]{babel}
\fi
\ifLuaTeX
  \usepackage{selnolig} % disable illegal ligatures
\fi


\setlength{\emergencystretch}{3em} % prevent overfull lines

\providecommand{\tightlist}{%
\setlength{\itemsep}{0pt}\setlength{\parskip}{0pt}}






\makeatletter
\@ifpackageloaded{tcolorbox}{}{\usepackage[skins,breakable]{tcolorbox}}
\@ifpackageloaded{fontawesome5}{}{\usepackage{fontawesome5}}
\definecolor{quarto-callout-color}{HTML}{909090}
\definecolor{quarto-callout-note-color}{HTML}{0758E5}
\definecolor{quarto-callout-important-color}{HTML}{CC1914}
\definecolor{quarto-callout-warning-color}{HTML}{EB9113}
\definecolor{quarto-callout-tip-color}{HTML}{00A047}
\definecolor{quarto-callout-caution-color}{HTML}{FC5300}
\definecolor{quarto-callout-color-frame}{HTML}{acacac}
\definecolor{quarto-callout-note-color-frame}{HTML}{4582ec}
\definecolor{quarto-callout-important-color-frame}{HTML}{d9534f}
\definecolor{quarto-callout-warning-color-frame}{HTML}{f0ad4e}
\definecolor{quarto-callout-tip-color-frame}{HTML}{02b875}
\definecolor{quarto-callout-caution-color-frame}{HTML}{fd7e14}
\makeatother
\makeatletter
\@ifpackageloaded{caption}{}{\usepackage{caption}}
\AtBeginDocument{%
\ifdefined\contentsname
  \renewcommand*\contentsname{Table des matières}
\else
  \newcommand\contentsname{Table des matières}
\fi
\ifdefined\listfigurename
  \renewcommand*\listfigurename{Liste des Figures}
\else
  \newcommand\listfigurename{Liste des Figures}
\fi
\ifdefined\listtablename
  \renewcommand*\listtablename{Liste des Tables}
\else
  \newcommand\listtablename{Liste des Tables}
\fi
\ifdefined\figurename
  \renewcommand*\figurename{Figure}
\else
  \newcommand\figurename{Figure}
\fi
\ifdefined\tablename
  \renewcommand*\tablename{Table}
\else
  \newcommand\tablename{Table}
\fi
}
\@ifpackageloaded{float}{}{\usepackage{float}}
\floatstyle{ruled}
\@ifundefined{c@chapter}{\newfloat{codelisting}{h}{lop}}{\newfloat{codelisting}{h}{lop}[chapter]}
\floatname{codelisting}{Listing}
\newcommand*\listoflistings{\listof{codelisting}{Liste des Listings}}
\makeatother
\makeatletter
\makeatother
\makeatletter
\@ifpackageloaded{caption}{}{\usepackage{caption}}
\@ifpackageloaded{subcaption}{}{\usepackage{subcaption}}
\makeatother
\usepackage{bookmark}
\IfFileExists{xurl.sty}{\usepackage{xurl}}{} % add URL line breaks if available
\urlstyle{same}
\hypersetup{
  pdftitle={CleanMyMap x World Cleanup Day France},
  pdflang={fr},
  colorlinks=true,
  linkcolor={blue},
  filecolor={Maroon},
  citecolor={blue},
  urlcolor={blue},
  pdfcreator={LaTeX via pandoc}}


\title{CleanMyMap x World Cleanup Day France}
\usepackage{etoolbox}
\makeatletter
\providecommand{\subtitle}[1]{% add subtitle to \maketitle
  \apptocmd{\@title}{\par {\large #1 \par}}{}{}
}
\makeatother
\subtitle{Proposition de pilote Île-de-France pour mieux suivre les
cleanwalks locales}
\author{}
\date{}
\begin{document}
\maketitle


\setstretch{1.08}
\begin{tcolorbox}[enhanced jigsaw, arc=.35mm, bottomrule=.15mm, bottomtitle=1mm, breakable, colback=white, colbacktitle=quarto-callout-note-color!10!white, colframe=quarto-callout-note-color-frame, coltitle=black, left=2mm, leftrule=.75mm, opacityback=0, opacitybacktitle=0.6, rightrule=.15mm, title=\textcolor{quarto-callout-note-color}{\faInfo}\hspace{0.5em}{Note}, titlerule=0mm, toprule=.15mm, toptitle=1mm]

Ce document va à l'essentiel. Il montre ce que CleanMyMap peut apporter
à World Cleanup Day France sur un pilote Île-de-France.

\end{tcolorbox}

\subsection{Ce que CleanMyMap propose}\label{ce-que-cleanmymap-propose}

CleanMyMap est un site open-source pensé pour donner une lecture simple
et exploitable des cleanwalks locales. L'outil relie la déclaration
d'action, le tracé GPS réel, la cartographie, les rapports d'impact et
les espaces d'échange.

\begin{itemize}
\tightlist
\item
  Suivre précisément les parcours réalisés.
\item
  Faire remonter les zones réellement couvertes.
\item
  Produire un rapport PDF clair après action.
\item
  Structurer un pilote local réutilisable.
\end{itemize}

\begin{tcolorbox}[enhanced jigsaw, arc=.35mm, bottomrule=.15mm, bottomtitle=1mm, breakable, colback=white, colbacktitle=quarto-callout-tip-color!10!white, colframe=quarto-callout-tip-color-frame, coltitle=black, left=2mm, leftrule=.75mm, opacityback=0, opacitybacktitle=0.6, rightrule=.15mm, title=\textcolor{quarto-callout-tip-color}{\faLightbulb}\hspace{0.5em}{Astuce}, titlerule=0mm, toprule=.15mm, toptitle=1mm]

L'objectif n'est pas de remplacer votre organisation, mais d'ajouter une
couche de lecture, de preuve et de coordination plus fine.

\end{tcolorbox}

\subsection{Ce que cela change pour WCD
France}\label{ce-que-cela-change-pour-wcd-france}

{\def\LTcaptype{none} % do not increment counter
\begin{longtable}[]{@{}
  >{\raggedright\arraybackslash}p{(\linewidth - 4\tabcolsep) * \real{0.3333}}
  >{\raggedright\arraybackslash}p{(\linewidth - 4\tabcolsep) * \real{0.3333}}
  >{\raggedright\arraybackslash}p{(\linewidth - 4\tabcolsep) * \real{0.3333}}@{}}
\toprule\noalign{}
\begin{minipage}[b]{\linewidth}\raggedright
Besoin côté WCD France
\end{minipage} & \begin{minipage}[b]{\linewidth}\raggedright
Ce que CleanMyMap apporte
\end{minipage} & \begin{minipage}[b]{\linewidth}\raggedright
Effet concret
\end{minipage} \\
\midrule\noalign{}
\endhead
\bottomrule\noalign{}
\endlastfoot
Mieux suivre les cleanwalks locales & tracé GPS réel des actions &
visualiser les parcours au lieu d'avoir seulement une déclaration
textuelle \\
Clarifier la saisie terrain & formulaire enrichi pour l'action & capter
poids, mégots, zones, photos et contexte sans alourdir le terrain \\
Lire l'activité d'un territoire & carte des actions et filtres & repérer
les zones nettoyées, les doublons et les secteurs peu couverts \\
Partager un bilan utile & rapport PDF automatique & transmettre un
document lisible à l'équipe, aux élus ou aux partenaires \\
Tester un usage régional & pilote Île-de-France & valider la méthode sur
un périmètre concret avant toute extension \\
Mieux coordonner les acteurs & espaces d'échange et pilotage & faciliter
la remontée d'information entre terrain, associations et coordination \\
\end{longtable}
}

\begin{tcolorbox}[enhanced jigsaw, arc=.35mm, bottomrule=.15mm, bottomtitle=1mm, breakable, colback=white, colbacktitle=quarto-callout-note-color!10!white, colframe=quarto-callout-note-color-frame, coltitle=black, left=2mm, leftrule=.75mm, opacityback=0, opacitybacktitle=0.6, rightrule=.15mm, title=\textcolor{quarto-callout-note-color}{\faInfo}\hspace{0.5em}{Note}, titlerule=0mm, toprule=.15mm, toptitle=1mm]

Lecture rapide: \texttt{tracé\ GPS} -\textgreater{}
\texttt{formulaire\ enrichi} -\textgreater{}
\texttt{carte\ des\ actions} -\textgreater{} \texttt{rapport\ PDF}
-\textgreater{} \texttt{pilote\ Île-de-France}

\end{tcolorbox}

\subsection{Les logiques métier}\label{les-logiques-muxe9tier}

Le cœur du projet tient en quelques logiques simples.

\begin{itemize}
\tightlist
\item
  Déclarer une action en quelques minutes, sans alourdir le terrain.
\item
  Vérifier qu'une action correspond bien à une zone réellement
  parcourue.
\item
  Comparer les opérations par association, ville, période ou type
  d'action.
\item
  Valoriser les initiatives avec des rapports plus lisibles pour les
  élus et partenaires.
\item
  Garder une séparation claire entre terrain, coordination et pilotage.
\end{itemize}

\begin{tcolorbox}[enhanced jigsaw, arc=.35mm, bottomrule=.15mm, bottomtitle=1mm, breakable, colback=white, colbacktitle=quarto-callout-tip-color!10!white, colframe=quarto-callout-tip-color-frame, coltitle=black, left=2mm, leftrule=.75mm, opacityback=0, opacitybacktitle=0.6, rightrule=.15mm, title=\textcolor{quarto-callout-tip-color}{\faLightbulb}\hspace{0.5em}{Astuce}, titlerule=0mm, toprule=.15mm, toptitle=1mm]

Pour le partenariat, l'intérêt principal est la combinaison entre
visibilité locale et lecture consolidée.

\end{tcolorbox}

\subsection{Le socle technique}\label{le-socle-technique}

{\def\LTcaptype{none} % do not increment counter
\begin{longtable}[]{@{}
  >{\raggedright\arraybackslash}p{(\linewidth - 2\tabcolsep) * \real{0.5000}}
  >{\raggedright\arraybackslash}p{(\linewidth - 2\tabcolsep) * \real{0.5000}}@{}}
\toprule\noalign{}
\begin{minipage}[b]{\linewidth}\raggedright
Brique
\end{minipage} & \begin{minipage}[b]{\linewidth}\raggedright
Rôle
\end{minipage} \\
\midrule\noalign{}
\endhead
\bottomrule\noalign{}
\endlastfoot
\texttt{Supabase} & base de vérité pour les données, les profils, les
actions et les règles \texttt{RLS} \\
\texttt{Quarto} & moteur de rendu du PDF à partir du Markdown \\
\texttt{GitHub} & dépôt source, historique des changements et
collaboration technique \\
\end{longtable}
}

Ce socle reste simple à comprendre pour un développeur: une base de
données unique, un rendu documentaire reproductible, et un historique de
code auditable. Le projet peut donc évoluer sans perdre la traçabilité
des choix.

\subsection{Suite possible}\label{suite-possible}

Si vous souhaitez aller plus loin, je peux préparer l'une des trois
versions suivantes :

\begin{itemize}
\tightlist
\item
  une démonstration terrain centrée sur le tracé GPS et le formulaire
  d'action ;
\item
  une version plus institutionnelle pour vos équipes ;
\item
  une proposition d'échange technique avec votre développeur autour de
  \texttt{Supabase}, des \texttt{RLS} et des flux de données.
\end{itemize}

Maxence Deroome\\
CleanMyMap\\
cleanmymap.fr




\end{document}
